\section{Summary}

Deep learning has advanced hydrological prediction, but model interpretability remains challenged by low verifiability, neglected variable dependencies\cite{e24050687_Shap_problem}, and unclear internal structures\cite{kumar2020_Shap_problem}, leaving uncertain whether high-performance models learn physically meaningful hydrological processes\cite{ReviewOfXAI}.
This study proposes the Sparse Transformer-based Explainable Hydrology (STEH) framework. It establishes an end-to-end interpretability pipeline using structured sparsification, circuit-level ablation, variable interaction analysis\cite{wang2022interpretabilitywildcircuitindirect}, and representation probing to systematically reveal the internal decision-making mechanisms of deep learning models in hydrology.

The main findings are as follows.First, STEH provides a circuit-level view to explicitly identify decision pathways and information flow, offering a structurally grounded and verifiable interpretability paradigm.Second, simpler models form more compact internal circuits, while over-complex models introduce redundant structures; model parsimony improves clarity while preserving physical meaning.Third, physically meaningful representations are learned across models regardless of performance, indicating physical knowledge is inherent in data-driven models.Fourth, STEH reliably characterizes stable and hydrologically consistent variable interactions across diverse hydro-climatic regimes.Fifth, sensitivity analysis shows the model effectively captures dominant streamflow-generation knowledge but incompletely represents temporal organization and state evolution.Sixth, STEH-derived knowledge is transferable and practically useful, as validated by performance improvements in a Random Forest model.Seventh, this study provides a replicable pathway for interpretability and physical validation of deep learning models in hydrological research.

In summary, this study demonstrates that hydrological deep learning models encode physically meaningful information that can be systematically uncovered through structural analysis. The STEH framework bridges black-box prediction and hydrological interpretation, supports knowledge discovery, and offers a generalizable route for integrating data-driven and physics-based hydrological modeling.